\section{讨论与相关工作}
\subsection{与校准和实验设计研究的关系}
\paragraph{不可识别情境下的校准。}
贝叶斯计算机模型校准区分参数不确定性与模型偏差 \citep{kennedy2001}。已有资格集合校准研究处理了多智能体模拟器的不可识别问题，并结合分布比较与贝叶斯优化 \citep{bai2022}。本文的经验过程／轨迹筛选更简单，不具有这些方法的统计检验解释。用于无似然推断的贝叶斯优化分配的是昂贵模拟器的评估次数 \citep{gutmann2016}，而本文在推演已经缓存后，分配的是过程记录的访问预算；两者对应不同的资源决策。因此，集合校准本身不是新颖性主张。本文着重把它与阶段取证机会、组合目标轨迹明确连接，并通过可审计实验研究这种连接何时失效。

\paragraph{预测导向取证。}
基于方差的主动学习已有长期研究 \citep{cohn1996}，贝叶斯实验设计则为取决于效用的数据获取提供了决策理论框架 \citep{chaloner1995}。目标导向 OED 关心预测量，而非全部参数不确定性 \citep{attia2018,zhong2026}。本文的协方差平方评分是有限机制、二元观测下的特例，采用方差效用而非预测信息增益。若各机制给某记录相同的概率，或该记录与探针预测不相关，即使行为熵很高，其取证价值也可能为零。实际关键在于阶段似然是否有效、探针目标分布是否合适、数值支持是否保留，以及计算能否复用。当前证据对计算复用的支持最明确。

\paragraph{社会模拟系统。}
CALM 将校准后的随机选择用于基于活动的出行模拟 \citep{cheng2026}，针对的是具体真实移动场景，而非本文的有限过程记录获取问题。SocioVerse2 为社会系统建立纵向模拟与研究循环 \citep{zhang2026}；Large Population Models 强调群体规模的行为与互动 \citep{chopra2025}。这些系统支持在机制、观测和集体结果之间建立可复用接口的研究方向，但并非本文已经运行的实证基线。给适配器起一个基准名称，也不等于复现了该基准的评测。提案中从宏观观测学习个体动态的更广泛目标，在 macro2mind 中也有所体现 \citep{macro2mind2026}；本文没有实现或比较这一模型。真实领域比较仍需要相容的观测机会、有授权的数据，以及具名且冻结的评分协议。

\subsection{适用范围、边界条件与失效模式}
\label{sec:boundaries}
\paragraph{权重集中会消除取证所依赖的信号。}
当 $w$ 趋近单点质量时，对每条记录均有 $\Cov_w(p_r,Z)\to0$，因而 $A_r(w)\to0$。所以，较大的硬相容集合不足以保证保留认识不确定性。较小核宽与有限精度 softmax 共同限制了有效支持；后验更新无法恢复恰好为零的候选质量。实测未发现取证收益，与这一极限行为一致，但要定量证明其贡献，还需超出本冻结研究的受控干预。

\paragraph{为什么一些看似合理的改动失败了。}
R8--R9 和 R36--R37 的较宽核被回退：在当前开发目标下，给更多候选保留质量增加了预测误差。R14--R17、R50--R53 混入均匀权重，也没有改善目标。这些结果区分了数值多样性与准确不确定性，并不证明多样化集合通常更差。R10--R11 将温度降到 $0.05$ 和 $0.1$、加强证据似然后，结果变差，这与有限机制库失配时少量噪声二元证据的过度影响相符。混合分布已经高度集中后，后续温度变化只带来极小差别。

R18--R23 加入过程 Brier 能量，所选行为几乎未变；R24--R29 调整通道权重，没有取得超过接受阈值的改善。R31--R35 的最差通道、绝对误差、Huber、截尾与后半程损失也未提高分数。这些损失体现不同的拟合观念，没有哪一个是普遍有效的稳健修正。R42--R45 重缩放探针坐标后 AUC 完全相同，因此不能把这些已回退的构想写成最终方法的创新。R54--R57 的种子边缘化核表现更差，但同时改变了校准构造与核尺度，不能据此断言随机性建模有害。最后，R59--R60 提高方差底值，使 R58 的收益消失；这种正则敏感性是一项限制，而不是所选参数能迁移到新领域的证据。

\paragraph{批量条件化不蕴含次模性。}
自适应次模保证需要边际收益递减性质 \citep{golovin2011}，本文目标一般不满足它。考虑以两个独立公平二元变量 $(U,V)$ 表示机制，目标 $Z=U\oplus V$ 是两者的异或，两个无噪声查询分别揭示 $U$ 和 $V$。任一查询单独出现时，目标方差都不变；两个一起则完全确定 $Z$。即使给定机制后查询是确定的、因而条件独立，第二个查询的边际价值仍在第一个查询后增加，违反收益递减。这个解析反例说明，精确假设分支期望应称为贪心条件化取证，不能赋予通用近似保证。

\paragraph{预测一致与任意程度的模型失配可以同时存在。}
假设所有候选预测 $z_0$，而一个遗漏的环境规则使现实变成 $z_0+\Delta$。所有认识协方差与取证评分都为零，但均方误差为 $\|\Delta\|^2/D$。方差缩减恒等式在模型内部依然成立，却无法发现候选类中根本不存在的机制；较宽的随机区间还可能掩盖弱偏移。事后残差检查提供另一种证据来源，但本研究中的阈值既没有校准的误报控制，也没有事前预警保证。扩充候选与显式偏差建模是可能的改进方向，而非已经实现的成功结果。

\paragraph{当前证据支持什么，不支持什么。}
同一个开发划分用于选择 60 个候选，只报告最优分数而不披露搜索历史，会夸大验证力度。测试评测与独立世界扩展增强了可复现性，但仍固定活动模板、训练数据、证据和小型有限机制库。文中的任何置信区间都没有覆盖真实人群不确定性或搜索过程本身的不确定性。第~\ref{sec:round2} 节使用新的查询群体和封存测试集，部分回应了这一问题。后续仍需改变事件设计、历史数据和证据池，加入同先验且优化实现的 OED，检验非坍缩取证收益能否超出本次构造的压力场景。还应把参数共享与先验变化分开，在使用真实行为记录前验证观测模型。

\subsection{结论}
本文明确并实现了一个计算成本较低的预算化过程取证闭环，支持完整轨迹预测和分解的不确定性。最重要的经验发现是：成功的群体校准可能消除主动取证赖以发挥作用的机制分歧。冻结实现改善了合成轨迹预测，并加速了参考评分计算，但 V1 的匹配对照没有建立主动取证或保留集合的优势。独立的 V2 压力测试展示了相对同先验 Random 和 Entropy 的取证收益，同时与 ParamIG 持平；其刻意设置的无关高熵记录和未生效的批量消融限定了结论范围。将这种区别如实作为结果，比未经支持地宣称机制已被识别或超越外部基准，更有助于为组合式社会模拟建立可靠研究基础。
